\section{Which small weights can occur?}
\label{sec:support}
Write $T_n=(\F^\times)^n$, and denote the number of nonzero coordinates of a vector $z$ by $\wt(z)$. The graph
\[
 C_A=\{(x,Ax):x\in\F^n\}
 \quad\text{has orthogonal complement}\quad
 C_A^\perp=\{(-A^{\mathsf T}y,y):y\in\F^n\}.
\]
Both coordinate projections on each code are bijective. We call $(\wt(x),\wt(z))$ the split weight of $(x,z)$. We call $x$ a solution for $A$ if $x$ and $Ax$ both lie in $T_n$. Such a solution is precisely the input of a word of $C_A$ of split weight $(n,n)$.

\begin{lemma}[Support exclusions]\label{lem:support}
Suppose that every invertible matrix of dimension less than $n$ over $\F$ has a solution. If $A\in\GL_n(\F)$ has none, then
\begin{enumerate}
\item $C_A$ contains no word of split weight $(i,j)$ with $1\le i,j\le3$.
\item $C_A^\perp$ contains no word of split weight $(1,1),(1,2),(2,1),(1,3),(3,1)$ or $(2,2)$.
\end{enumerate}
\end{lemma}

\begin{proof}
For the first assertion, regard the standard coordinate vectors and the rows of $A$ as two bases of a coefficient space $E$. A functional nonzero on both bases gives the required point by taking its values on the first basis. In any space of dimension less than $n$, the induction hypothesis lets us find such a functional even if we prescribe only two independent families: extend each family to a basis and use the change-of-basis matrix.

Suppose a nonzero functional $d$ has nonzero values on precisely $i$ members of the first basis and $j$ members of the second. These are the split weights of the word defined by $d$. Call these basis vectors moving and the others fixed. If $i+j\le4$, apply induction inside $\ker d$ to the two fixed families, then extend the resulting functional $L$ to $E$. Each moving vector $v$ forbids just one value of $c$ for $L+cd$, namely $-L(v)/d(v)$. At most four values are forbidden in a field of size five.

It remains to consider $i,j\in\{2,3\}$ with $i+j\ge5$. Let $U,V$ be the respective spans of the fixed vectors, of dimensions $n-i,n-j$. Let $P$ be the set of normalized moving vectors $v/d(v)$ from either basis. Then $U+\langle P\rangle=V+\langle P\rangle=E$. We shall identify at least $k=\min(i,j)-1$ pairs successively, without introducing a dependence in either fixed family. Starting with $S=0$, suppose $\dim S=r<k$ and $S$ meets each of $U,V$ trivially. The images of $P$ in $E/(U+S)$ cannot all be equal: otherwise $E$ would be spanned by $U+S$ and one more vector, although $n-i+r+1<n$. The same holds in $E/(V+S)$. Two nonconstant maps on a common set have a pair on which both change. Indeed, if the second map takes equal values on every pair separated by the first, comparison with a point in a different fiber shows that the second is constant. Adjoin to $S$ the difference of such a pair of points of $P$. Its two images are nonzero, so the new subspace is still disjoint from each fixed span. This step strictly reduces the number of distinct images of $P$ modulo $S$.

After $k$ steps the two fixed families remain independent in $E/S$. Choose a functional there by induction and pull it back to $L$ on $E$. Now at most $|P|-k\le i+j-k=\max(i,j)+1\le4$ distinct values of $-L(v/d(v))$ occur. Again some $L+cd$ avoids every zero. This proves the first assertion, including the pair $(3,3)$ used in dimension eleven.

For the second assertion, a dual word with right weight one and left weight at most three identifies a row $R$ supported on at most three coordinates. Write $w=A^{-1}e_s$, where $s$ is its row index. Delete output $s$ after restricting inputs to a hyperplane of the form $x_j=c x_k$, $c\ne0$. The induced square map is invertible whenever $w_j\ne c w_k$: a kernel vector would map to a multiple of $e_s$ and hence would be a multiple of $w$. Any nowhere-zero solution for the smaller map has all original input coordinates nonzero. It therefore suffices to choose the hyperplane so row $R$ is nonzero there on every such input.

For a row on at most two sites, take $j,k$ to include its support with $R_j\ne0$ (the one-dimensional case is immediate). We require $cR_j+R_k\ne0$ and $w_j\ne c w_k$. The latter is not the zero condition identically, since $Rw=1$. At most two of the four nonzero scalars are excluded. For a row with three terms, at least one pair $j,k$ has $R_j w_j+R_k w_k\ne0$, since the sum over all three pairs equals $2Rw=2$. Set $c=-R_k/R_j$. This hyperplane misses $w$ and cancels exactly the chosen two terms, leaving a nonzero multiple of the third coordinate. Applying the same argument to $A^{-1}$ exchanges the two blocks and covers the reverse weight pairs.

Finally, a dual word $(u,v)$ of split weight $(2,2)$ gives $u=-A^{\mathsf T}v$. The map $A$ sends $\ker u^{\mathsf T}$ isomorphically onto $\ker v^{\mathsf T}$. On each hyperplane two coordinates are nonzero scalar multiples of one another. Deleting one of them on each side identifies the restriction with an invertible $(n-1)$-dimensional map; reinserting it preserves nonvanishing on both sides. This proves the second assertion.
\end{proof}

\section{Ordinary weights and the ten-dimensional theorem}
\label{sec:ordinary}
Let $H_{ij}$ and $D_{ij}$ count the words of split weight $(i,j)$ in $C_A$ and $C_A^\perp$, respectively. Their row and column sums are the integers $\mu_i=\binom ni4^i$, indexed by the weight in the relevant block. We use the symmetric tables $W=H+H^{\mathsf T}$ and $V=D+D^{\mathsf T}$. They have origin entry $m=2$ and row sums $m\mu_i$. In particular $W_{nn}=2H_{nn}$; this symmetrization preserves the question we want to answer.

For $0\le a,i\le n$ put
\[
 K_{ai}=[z^a](1+4z)^{n-i}(1-z)^i.
\]
If $\chi$ is a nontrivial additive character of $\F$, then for a fixed word $x$ of weight $i$ the sum of $\chi(u\cdot x)$ over $\wt(u)=a$ is $K_{ai}$. A zero coordinate of $x$ contributes four characters of value one at degree one; at a nonzero coordinate their sum is $-1$. Summing the character identity
$\sum_{c\in C_A}\chi(u\cdot c)=5^n{\bf1}_{u\in C_A^\perp}$
over an ambient split-weight class, and then adding the transposed identity, gives
\begin{equation}
 5^n V=K W K^{\mathsf T}.
 \label{eq:ordinary-transform}
\end{equation}
MacWilliams showed how the ordinary weight distribution of a binary linear code determines the distribution of its dual \cite{macwilliams}. Here the same character argument keeps the two block weights separate, which is what allows us to single out the cell $(n,n)$.

\begin{proposition}[An integer certificate]\label{prop:ordinary}
Let $W,V$ be symmetric nonnegative tables with $W_{00}=V_{00}=m>0$ and row sums $m\mu_i$ in each table. In dimensions $6\le n\le10$, \eqref{eq:ordinary-transform} and the support exclusions of Lemma~\ref{lem:support} imply $W_{nn}>0$. In dimension $11$ the same conclusion holds under the additional assumption $V_{23}=0$.
\end{proposition}
\begin{proof}
Choose the two integer vectors $u,v$ and the integer matrix $T$ specified for $n$ in the accompanying certificate. Set
\begin{align*}
 P_{ij}&=u_i+u_j+(K^{\mathsf T}TK)_{ij},&
 Q_{ij}&=v_i+v_j-5^nT_{ij},\\
 R&=2\sum_i (u_i+v_i)\mu_i .
\end{align*}
With $\langle X,Y\rangle=\sum_{ij}X_{ij}Y_{ij}$, pairing \eqref{eq:ordinary-transform} with $T$ cancels the two transform terms. The marginals give the exact identity
\[
 \langle W,P\rangle+\langle V,Q\rangle=mR.
\]
Suppose $W_{nn}=0$. Apart from the two origins, all coefficients $P_{ij},Q_{ij}$ are nonnegative on cells permitted to carry mass. In dimension eleven the test for $Q$ also exempts $(2,3),(3,2)$; they have zero mass by the extra hypothesis and symmetry. Consequently the left side is at least $m(P_{00}+Q_{00})$. Table~\ref{tab:gaps} gives the strict difference per unit mass. Each positive difference contradicts the identity.
\end{proof}

\begin{table}[ht]
\centering
\caption{Integer gaps $P_{00}+Q_{00}-R$. Only the eleven-dimensional entry requires $V_{23}=0$.}\label{tab:gaps}
\begin{tabular}{rr@{\qquad}rr}\toprule
$n$ & Gap & $n$ & Gap \\\midrule
6 & 7812500000 & 9 & 286749531250000000000\\
7 & 5860862993750000000 & 10 & 235472537548828125000000\\
8 & 1256583418475000000000 & 11 & 170253422831819243164062500000\\
\bottomrule\end{tabular}
\end{table}

For clarity, the ordinary certificate uses only a subset of Lemma~\ref{lem:support}. The primal sign test exempts the origin, $(n,n)$, and $1\le\min(i,j)\le2$, $\max(i,j)\le3$. The dual test exempts the origin and the pairs $\min(i,j)=1$, $\max(i,j)\le2$, together with the extra pair and its transpose for $n=11$. All signs and all six gaps can be checked from the integer arrays by the short replay in Appendix~\ref{sec:replay}. They require no optimization or rounding.

\begin{corollary}\label{cor:ten}
Every $A\in\GL_n(\F)$ with $n\le10$ has a solution. If an eleven-dimensional counterexample exists, its dual code contains a word with block weights $(2,3)$ or $(3,2)$.
\end{corollary}
\begin{proof}
For $n\le5$ the assertion follows from Yu's bound \cite{yu}. In a least failing dimension between six and ten, Lemma~\ref{lem:support} applies. Proposition~\ref{prop:ordinary} now gives $H_{nn}>0$, a contradiction. In dimension eleven the proposition implies $V_{23}>0$ in a counterexample. Since $V_{23}=D_{23}+D_{32}$, at least one of these counts in the actual dual code is positive.
\end{proof}
